\scanpage{180}
\mspage{96}
% archive sheet number: 179
% lines 1-3
a) (pour mémoire) le foncteur canonique
\[
C\longrightarrow\operatorname{Ind}_{\pi}(C)\unsure{\xrightarrow{\approx}|\approxeq|\simeq|\cong}\operatorname{Ind}'_{\pi}(C)
\]
commute aux $\unsure{\underline{\lim}|\varinjlim}$ de cardinal $\leq\pi$, et

% lines 4-11
b) (pour mémoire) les systèmes inductifs de cardinal $\leq\pi$ dans $\operatorname{Ind}_{\pi}(C)$ (i.e. les foncteurs $\varphi:J\mathrel{\rule{1.5em}{0.4pt}}\operatorname{Ind}_{\pi}(C)$, avec $\operatorname{card}\operatorname{Fl}J\leq\pi$) admettent une représentation indicielle
\[
\psi:I\times J\longrightarrow C,
\]
avec $I$ ens. ordonné grand devant $\pi$, et admettent une limite inductive dans $\operatorname{Ind}_{\pi}C$ qui se calcule par
\[
\underbrace{\varinjlim_{j\in J}\varphi(j)}_{\text{dans }\operatorname{Ind}'_{\pi}(C)}\simeq\text{``}\varinjlim_{i\in I}\text{''}\Psi(i)
\]
où
\[
\Psi(i)=\varinjlim_{j\in J}{}_{C}\psi(i,j).
\]

% lines 17-18
c) $\operatorname{Ind}_{\pi}C$ est stable par petites limites inductives.

% margin note written sideways
\margin{\unsure{rajouter} ici, en d), le cor. 4.4.\unsure{16|15} \unsure{pge|page} 101}

% bottom paragraph
Il reste à établir ce dernier point. Comme on a déjà les $\underline{\lim}$ finies par b), il reste à prouver l'existence des \inserted{petites} $\underline{\lim}$ filtrantes, \deleted{\unsure{des} $\underline{\lim}$ $\unsure{I}$ \unsure{quelconques}} \inserted{\unsure{ou} \unsure{en} \illegible{1}} \deleted{filtrantes} des sommes infinies, ce qui est plus facile en fait. Mais
